\documentclass[12pt]{article}
\usepackage[utf8]{inputenc}
\usepackage[spanish]{babel}
\usepackage{graphicx}
\usepackage{geometry}
\geometry{margin=2.5cm}

\title{\Huge Bitácora 1\\ \Large Análisis de la relación entre el uso de redes sociales y el rendimiento académico en estudiantes universitarios}
\author{Katia Moreno Aspiroz C6A060\\
Jeremy Garcia Solano C33171}
\date{\today}

\begin{document}
\maketitle
\newpage

\tableofcontents
\newpage

\section{Introducción}
\subsection{Contexto}
El uso de redes sociales forma parte del día a día de la mayoría de universitarios. Plataformas como TikTok, Instragram o WhatsApp ocupan un tiempo significativo a lo largo día, incluso en momentos dedicados al estudio o a pausas durante el mismo.
\\
En este contexto, resulta interesante analizar si existe una relación entre el tiempo dedicado a las redes sociales y el rendimiento académico, ya que es razonable pensar que estos hábitos pueden influir en la concentración, y en consecuencia, en los resultados académicos.

\subsection{Motivación}
La motivación de este proyecto surge

\subsection{Audiencia}

\subsection{Posicionamiento}


\section{Idea}
\subsection{Definición}
\subsection{Conceptualización}
\subsection{Identificación de tensiones}
\subsection{Reformulación en modo pregunta}
\subsubsection{Preguntas propuestas}
\subsubsection{Pregunta seleccionada}
\subsubsection{Justificación de descarte}
\subsection{Argumentación de la pregunta}


\section{Datos y argumentación}
\subsection{Fuente de información}
Para este proyecto se utiliza una base de datos obtenida del sitio web Kaggle, la cual según su autor, fue elaborada para elaborar una tesis de pregrado titulada "The Impact of Social Media on Academic and Personal Well-being: A Data-Driven Approach".
\subsection{Contexto temporal y espacial}
Los datos recolectados en la base de datos, incluye principalmente estudiantes de la "Metropolitan University" y otras comunidades educativas que no se especifican en la descripción de la base de datos.
\subsection{Población y muestra}
La Población la cual se estudia son los estudiantes universitarios de diferentes comuninades académicas.
La muestra corresponde a las 404 lineas que están presentes en el la base de datos, donde cada linea representa a un estudiante que realizó la encuesta. 
\subsection{Unidad estadística}
Cada estudiante que realizó la encuesta, corresponde a la unidad estadística.

\end{document}
